\scanpage{18}
an algebro-geometric interpretation in terms of \inserted{moduli for} motives, is the following:

start with a vector space of finite \unsure{:|{}}\unsure{dimension-|dimension} $V$ over $\mathbb{Q}$, with \unsure{an|a} fundamental bilinear form $\unsure{\varphi|\phi}$ given, symmetric or alternating, take $G=\mathrm{SO}(\varphi)$, take any $i_0:U_{\mathbb{R}}\to G_{\mathbb{R}}$ s.th. $i_0(-1)=\mathrm{id}$ (if $\varphi$ symmetric), $i_0(-1)=-\unsure{\mathrm{Id}|\mathrm{id}}$ (if $\varphi$ \deleted{\unsure{a|o}} skew-symmetric), \deleted{$\unsure{i_0|\mathrm{id}}(\sqrt{-1})\unsure{=|\ne}\unsure{1|{-1}}$} \deleted{$i_0(i)$ \unsure{non} trivial.} \inserted{and a positivity condition of $\varphi$ relative to the bigrading of $V_{\mathbb{C}}$ defined by $i_0$, which I will not write down.}

Then $\unsure{\gamma|\mathscr{T}}$ is a direct generalization of Siegel's half-plane, and classifies all Hodge-structures \deleted{\unsure{under}} with underlying $\deleted{\unsure{V|Q}}\mathbb{Q}$-vector space $V$, \deleted{\unsure{(|a}} i.e. all bigradings of $V_{\mathbb{C}}$, \deleted{\unsure{having}} \inserted{satisfying} the usual symmetry condition for complex conjugation, which are compatible with $\varphi$ (\unsure{i.e|i.e.} $\varphi$ is a polarization), and for which the $\dim V_{\mathbb{C}}^{pq}$ have given values \unsure{.|{}} (depending on the weights of $U_{\deleted{\mathbb{R}}\mathbb{C}}$ operating on $V_{\mathbb{C}}\ \ldots$). The trouble is that, \deleted{\unsure{if|of}} as these special types of $\gamma$, $\gamma/\Gamma$ are in a way ``universal'' for all others
% archive sheet number: 17
